%% kapa: exam=1 round=36 year=2025 period=1 form=A subject=real_estate subject_ko=부동산학원론 seq=17 number=97
%% source: Q-Net 감정평가사 기출문제 (cst00309) artlSeq=5231525 file=raw/1008/5231525/2025년 제36회 감정평가사 1차시험 1교시 문제지 원본.pdf
%% conversion: kordoc 4.24.51 -> tools/gen1.py (auto); review=proofread-sonnet; flags=none
\begin{question}{97}
부동산의 투자과정에서 수익률에 관한 설명으로 옳은 것은? (단, 주어진 조건에 한함)
\begin{choices}
  \item 기대수익률은 본질적으로 사후수익률을 의미한다.
  \item 기대수익률은 시장이자율에 비례하고, 자산의 위험에도 비례한다.
  \item 기대수익률이 요구수익률보다 높으면, 대상부동산의 수요가 증가하여 요구수익률이 하락한다.
  \item 명목이자율로서 무위험이자율은 실질이자율에서 물가상승률을 차감한 값이다.
  \item 내부수익률이 요구수익률보다 큰 경우나 순현재가치가 1보다 큰 경우에는 투자하지 않는다.
\end{choices}
\end{question}
\answer{2}
